%**************************************************************************
%*  P01 -- Fair & Explainable Skin-Lesion Classification: CNN vs ViT       *
%*  ELCVIA camera-ready format (LaTeX2e)                                   *
%**************************************************************************

\documentclass[twoside,a4paper,11pt]{article}
\usepackage{times}
\usepackage{elcvia}
\usepackage[english]{babel}
\usepackage{graphicx}
\usepackage{amsmath}
\usepackage{booktabs}

% INFORMATION TO FILL IN CAMERA READY VERSION
%------------------------------------------------------------------
\newcommand{\vol}{0}            % Number of Volume (default 0)
\newcommand{\num}{0}            % Number (default 0)
\setcounter{page}{1}            % First page (default 1)
\newcommand{\begpag}{\thepage}
\newcommand{\finpag}{10}        % Last Page (update once final)
\newcommand{\yearp}{2026}       % Year of the issue
%------------------------------------------------------------------
\renewcommand{\thefootnote}{\fnsymbol{footnote}}

\title{Fair \& Explainable Skin-Lesion Classification: CNN vs Vision
Transformer
\footnote[0]{Correspondence to: $<$su-24013@sitare.org$>$ \vspace{0.2cm}}
\footnote[0]{\noindent Recommended for acceptance by $<$name$>$}
}

\author{Afzal Raza$^*$, Anamika Kumari$^*$, Khushi Narwariya$^*$ and Karanpal$^*$ \\ \\
\centerline{\small \em $^*$ Team 1, P01 -- HAM10000 / Fitzpatrick17k Fairness Study} \\ \\
\centerline{\small Received 1 January 2026; revised --; accepted --}
}

\begin{document}

\maketitle

\pagestyle{myheadings}

\markboth{\centerline{\small \it Raza et al. /
           Electronic Letters on Computer Vision and Image Analysis
           \vol(\num):\begpag-\finpag, \yearp}}
         {\centerline{\small \it Raza et al. /
           Electronic Letters on Computer Vision and Image Analysis
           \vol(\num):\begpag-\finpag, \yearp}}
\hrulefill

\begin{abstract}

We study whether a convolutional network (ResNet-50) or a Vision
Transformer (ViT-B/16) classifies the seven HAM10000 dermoscopic lesion
types more fairly across Fitzpatrick skin types (I--VI), and which model
yields more trustworthy visual explanations. Both architectures are
fine-tuned on HAM10000 with a lesion-level, leakage-free train/val/test
split, and evaluated for class-balanced accuracy, per-skin-tone
performance gaps (light I--III vs.\ dark IV--VI, using Fitzpatrick17k),
and explanation quality: Grad-CAM for the CNN and attention rollout for
the ViT, scored with deletion/insertion AUC. \emph{[Results to be
completed once training and evaluation are run; see
Section~\ref{sec:results}.]}

\vspace{0.5cm}
\noindent
{\em Key Words}: Computer Vision, Image Analysis, Fairness in Machine
Learning, Skin-Lesion Classification, Convolutional Neural Networks,
Vision Transformers, Explainability, Grad-CAM, Attention Rollout.

\end{abstract}
\hrulefill

\section{Introduction}
\label{sec:intro}

Automated dermoscopic image classification can help triage suspicious
skin lesions, but prior work has shown that deep models trained
predominantly on light-skinned patients can under-perform on darker skin
tones~\cite{Groh:21}. This raises both a scientific question --
\emph{which architecture generalizes more fairly?} -- and a practical
one -- \emph{how do we know what the model is actually looking at?}

This paper compares two widely used backbones, a CNN (ResNet-50) and a
Vision Transformer (ViT-B/16), fine-tuned on the HAM10000
dataset~\cite{Tschandl:18} for 7-class lesion classification
(\texttt{akiec}, \texttt{bcc}, \texttt{bkl}, \texttt{df}, \texttt{mel},
\texttt{nv}, \texttt{vasc}). We evaluate skin-tone fairness using the
Fitzpatrick17k dataset~\cite{Groh:21}, mapped onto the HAM10000 label
space, and compare the explanations each model produces -- Grad-CAM for
the CNN~\cite{Selvaraju:17} and attention rollout for the
ViT~\cite{Abnar:20,Dosovitskiy:21}.

\textbf{Research question:} Is ViT or CNN fairer across skin tones, and
which one gives better explanations?

\section{Related Work}
\label{sec:related}

Esteva et al.~\cite{Esteva:17} first showed dermatologist-level skin
cancer classification with CNNs. Groh et al.~\cite{Groh:21} introduced
Fitzpatrick17k and demonstrated accuracy disparities across skin tones
for clinical-image classifiers. Dosovitskiy et al.~\cite{Dosovitskiy:21}
introduced the Vision Transformer; Abnar and Zuidema~\cite{Abnar:20}
proposed attention rollout as a way to visualize what a transformer
attends to, analogous to Grad-CAM~\cite{Selvaraju:17} for CNNs. To our
knowledge, no prior work directly compares CNN vs.\ ViT fairness
\emph{and} explanation quality jointly on this pair of datasets.

\section{Datasets}
\label{sec:datasets}

\begin{description}
\item[HAM10000] 10{,}015 dermoscopic images across 7 diagnostic
categories, used for training, validation and testing. Splits are
constructed at the \emph{lesion} level (not image level) using a
stratified group split, so that no \texttt{lesion\_id} appears in more
than one split -- preventing leakage from images of the same lesion.

\item[Fitzpatrick17k] 16{,}577 clinical images labelled with Fitzpatrick
skin type (I--VI) and one of 114 clinical conditions. We map the subset
of conditions that correspond to our 7 HAM10000 classes
(\texttt{mel}, \texttt{bcc}, \texttt{akiec}, \texttt{bkl}, \texttt{df};
mappings for \texttt{nv} and \texttt{vasc} are still being finalized)
and use it purely as a skin-tone-labelled \emph{evaluation} set, since
its images are clinical rather than dermoscopic -- a domain shift we
report honestly rather than hide.
\end{description}

\section{Method}
\label{sec:method}

\subsection{Models}
Both backbones are initialized from ImageNet-pretrained weights and
fine-tuned with a new 7-class head. The CNN is a ResNet-50 (frozen
backbone for the first two epochs, then unfrozen); the ViT is
ViT-B/16 via \texttt{timm}, with layer-wise learning-rate decay to
preserve pretrained low-level features.

\subsection{Class-imbalance handling}
HAM10000 is heavily imbalanced towards the \texttt{nv} (melanocytic
nevus) class. We compare plain cross-entropy, class-weighted
cross-entropy, and focal loss~\cite{Lin:17} ($\gamma=2$), each with and
without a balanced \texttt{WeightedRandomSampler}.

\subsection{Explainability}
For the CNN we compute Grad-CAM~\cite{Selvaraju:17} heat-maps from the
last convolutional block. For the ViT we compute attention
rollout~\cite{Abnar:20} across all transformer blocks, projecting the
CLS-token attention back onto the $14\times14$ patch grid and
upsampling to the input resolution. Both explanation types are
normalized to $[0,1]$ heat-maps of the same shape for a like-for-like
comparison.

\subsection{Fairness and explanation-quality metrics}
We report overall and per-class balanced accuracy, macro-F1 and
macro-AUROC (Section~\ref{sec:method}), then, for Fitzpatrick17k,
balanced accuracy per Fitzpatrick type, the gap (max $-$ min) and ratio
(min~/~max) across types, and the light (I--III) vs.\ dark (IV--VI)
split, with bootstrap 95\% confidence intervals. Explanation quality is
scored with deletion and insertion AUC~\cite{Petsiuk:18}, compared
between Grad-CAM and attention rollout, overall and per skin-tone group.

\section{Experimental Setup}
\label{sec:setup}

All experiments use $224\times224$ inputs, ImageNet normalization,
AMP mixed-precision training, early stopping on validation balanced
accuracy, and a fixed seed (42) for reproducibility. Hyperparameters
(learning rate, batch size, epochs, schedule) are tuned for both models
under an equal compute budget; full configs are version-controlled in
\texttt{configs/resnet50.yaml} and \texttt{configs/vit\_b16.yaml}.

\section{Results}
\label{sec:results}

\emph{[Placeholder -- to be filled in once \texttt{scripts/train.py}
and \texttt{scripts/evaluate.py} are implemented and run. Planned
tables: (1) overall test accuracy / balanced accuracy / macro-F1 /
macro-AUROC for ResNet-50 vs.\ ViT-B/16; (2) per-Fitzpatrick-type
balanced accuracy and light-vs-dark gap for both models with bootstrap
CIs; (3) deletion/insertion AUC for Grad-CAM vs.\ attention rollout,
overall and per skin-tone group. See \texttt{results/EXPERIMENTS.md}
for the running experiment log.]}

\section{Discussion}
\label{sec:discussion}

\emph{[Placeholder -- interpretation of the fairness gap and
explanation-quality results above, once available, including whether
the more accurate model is also the fairer and more explainable one, or
whether these trade off against each other.]}

\section{Conclusion}
\label{sec:conclusion}

\emph{[Placeholder -- summary answer to the research question in
Section~\ref{sec:intro}, and limitations: domain shift between
HAM10000 (dermoscopic) and Fitzpatrick17k (clinical), and the small
number of dark-skin-tone samples available for evaluation.]}

\section*{Acknowledgements}

Team 1: Afzal Raza (data pipeline, splits, Fitzpatrick label mapping,
EDA), Anamika Kumari (ResNet-50 baseline, focal loss, training loop,
Grad-CAM), Khushi Narwariya (ViT-B/16 fine-tuning, attention rollout,
hyperparameter search), Karanpal (fairness and explanation-quality
evaluation, statistical tests).

\begin{thebibliography}{9}

\bibitem{Esteva:17}
A. Esteva et al., ``Dermatologist-level classification of skin cancer
with deep neural networks'', \emph{Nature} 542:115--118, 2017.

\bibitem{Dosovitskiy:21}
A. Dosovitskiy et al., ``An Image is Worth 16x16 Words: Transformers
for Image Recognition at Scale'', \emph{ICLR}, 2021.

\bibitem{Groh:21}
M. Groh et al., ``Evaluating Deep Neural Networks Trained on Clinical
Images in Dermatology with the Fitzpatrick 17k Dataset'',
\emph{CVPR Workshops}, 2021.

\bibitem{Tschandl:18}
P. Tschandl, C. Rosendahl, H. Kittler, ``The HAM10000 dataset'',
\emph{Scientific Data} 5:180161, 2018.

\bibitem{Selvaraju:17}
R. R. Selvaraju et al., ``Grad-CAM: Visual Explanations from Deep
Networks via Gradient-based Localization'', \emph{ICCV}, 2017.

\bibitem{Abnar:20}
S. Abnar, W. Zuidema, ``Quantifying Attention Flow in Transformers'',
\emph{ACL}, 2020.

\bibitem{Lin:17}
T.-Y. Lin et al., ``Focal Loss for Dense Object Detection'',
\emph{ICCV}, 2017.

\bibitem{Petsiuk:18}
V. Petsiuk, A. Das, K. Saenko, ``RISE: Randomized Input Sampling for
Explanation of Black-box Models'', \emph{BMVC}, 2018.

\end{thebibliography}

\end{document}
